% Part: first-order-logic
% Chapter: syntax-and-semantics
% Section: introduction

\documentclass[../../../include/open-logic-section]{subfiles}

\begin{document}

\olfileid{fol}{syn}{int}

\olsection{Pambuka}

Kanggo ngembangaké téori lan métatéori logika tataran kapisan,
kita kudu luwih dhisik netepaké sintaksis lan semantik
ekspresiné. Ekspresi logika tataran kapisan awujud suku lan
!!{formula}s. Suku dibangun saka !!{variable}s, !!{constant}s,
lan !!{function}s. Sabanjuré, !!^{formula}s dibangun saka
!!{predicate}s bebarengan karo suku (iki mbentuk !!{formula}s
paling prasaja, yaiku kang ``atomik''). Saka !!{formula}s atomik,
kita bisa mbangun wangun kang luwih rumit nganggo konektif
logika lan kuantor. Aturan pambentukan bisa ditetepaké kanthi
akèh cara; ing kéné kita mung mènèhi salah siji cara.
Sistem liya bisa milih lambang kang béda, milih konektif
primitif kang béda, lan nganggo kurung kanthi cara kang béda
(utawa ora nganggo kurung babar pisan, kaya notasi Polandia).
Nanging kabèh cara mau padha-padha netepaké himpunan suku lan
!!{formula}s sacara \emph{induktif}. Yèn aturan iki dirumusaké
kanthi trep, saben ekspresi dhasaré mung bisa kabangun kanthi
siji cara. Definisi induktif kang ngasilaké ekspresi kang
\emph{bisa diwaca kanthi unik} uga ngidini kita mènèhi teges
marang ekspresi mau kanthi cara kang padha: definisi induktif.

Nemtokaké teges ekspresi iku pakaryané semantik. Konsep
utama semantik yaiku relasi nyukupi ing !!a{structure}.
!!^a{structure} mènèhi teges marang unsur-unsur dhasar basa:
!!a{domain} iku himpunan obyek kang ora kosong. Kuantor
ditafsiraké ngliwati domain iki, !!{constant}s diwènèhi
unsur ing domain, !!{function}s diwènèhi fungsi saka
pangkat Kartesius domainé kang cocog karo aritasé menyang
!!{domain}, lan !!{predicate}s diwènèhi relasi ing
!!{domain}. !!{domain} bebarengan karo penetapan marang
kosakata dhasar mbentuk !!a{structure}. !!^{variable}s bisa
katon ing !!{formula}s; kanggo mènèhi semantik, kita uga
kudu netepaké !!{element}s saka !!{domain} kanggo variabel mau.
Iki diarani pangaji variabel. Pungkasané, relasi nyukupi
nyawijèkaké kabèh mau. !!^a{formula} bisa dicukupi ing
!!a{structure}~$\Struct{M}$ relatif marang pangaji variabel~$s$,
lan ditulis $\Sat{M}{!A}[s]$. Relasi iki uga ditetepaké
kanthi induksi tumrap wangun~$!A$. Tabel bebener kanggo
konektif logika, umpamané, netepaké kapan $!A \land !B$
dicukupi adhedhasar dicukupi (utawa ora) $!A$ lan ~$!B$.
Pangaji variabel banjur ora nduwèni pangaruh yèn
!!{formula}~$!A$ iku !!a{sentence}, yaiku ora nduwèni
variabel bebas. Mula kita bisa nyebut !!{sentence}s kang
dicukupi (utawa ora) ing !!{structure}s tanpa nyebut
pangaji variabel.

Adhedhasar relasi nyukupi $\Sat{M}{!A}$ kanggo ukara,
kita bisa netepaké konsep semantik dhasar: validitas,
akibat semantis, lan bisa dicukupi. Ukara diarani valid,
$\Entails !A$, yèn saben struktur nyukupiné. Ukara iku
dadi konsekuènsi saka himpunan !!{sentence}s,
$\Gamma \Entails !A$, yèn saben !!{structure} kang
nyukupi kabèh !!{sentence}s ing~$\Gamma$ uga nyukupi~$!A$.
Himpunan ukara bisa dicukupi yèn ana !!{structure} kang
nyukupi kabèh !!{sentence}s ing kono bebarengan. Amarga
!!{formula}s ditetepaké sacara induktif, lan relasi nyukupi
uga ditetepaké kanthi induksi tumrap wangun !!{formula}s,
kita bisa migunakaké induksi kanggo mbuktèkaké sipat-sipat
semantik lan sesambungan antarkonsep semantik kasebut.




\end{document}
